% !PRINTED 14
\weberpage{14}

Die Wurzeln der kubischen Funktion
\[
\varphi(z) = 4z^3 - g_2 z - y_3 % [sic] printed y_3; expected g_3
\]
mögen mit $e_1$, $e_2$, $e_3$ bezeichnet sein. Dann ist
\[
\varphi(z) = 4(z-e_1)(z-e_2)(z-e_3)
\]
und $e_1 + e_2 + e_3 = 0$.

Wir lassen die Werte von $x$ und $z$ einander folgendermaßen entsprechen:
\[
\begin{array}{ccccc}
x, & x_1, & x_2, & x_3, & x_4, \\
z, & \infty, & e_1, & e_2, & e_3,
\end{array}
\]
und es ergibt sich, wenn wir mit $m$ einen konstanten Faktor bezeichnen, der willkürlich angenommen werden kann:
\weq{4}{
\begin{aligned}
z - e_1 &= \frac{m(xx_2)}{(x_1x_2)(xx_1)} = m\left(\frac{1}{(x_1x_2)} - \frac{1}{(x_1x)}\right)\\
z - e_2 &= \frac{m(xx_3)}{(x_1x_3)(xx_1)} = m\left(\frac{1}{(x_1x_3)} - \frac{1}{(x_1x)}\right).\\
z - e_3 &= \frac{m(xx_4)}{(x_1x_4)(xx_1)} = m\left(\frac{1}{(x_1x_4)} - \frac{1}{(x_1x)}\right)
\end{aligned}
}

Der Faktor $m$ muß in allen drei Gleichungen derselbe sein, damit die Differenzen $(z-e_3)-(z-e_2)=e_2-e_3$ usw. von $x$ unabhängig werden.

Bildet man diese Differenzen und setzt wie in Algebra Bd.\,1, \S\ 70
\[
\begin{aligned}
(x_1x_2)(x_3x_4) &= U,\\
(x_1x_3)(x_4x_2) &= V,\\
(x_1x_4)(x_2x_3) &= W,
\end{aligned}
\]
so folgt
\[
e_2 - e_3 = \frac{-mU}{(x_1x_2)(x_1x_3)(x_1x_4)}.
\]

Führt man einen neuen willkürlichen Faktor $\mu$ ein, indem man
\weq{5}{m = 3\mu a_0 (x_1x_2)(x_1x_3)(x_1x_4)}
setzt, so folgt
\[
e_2 - e_3 = -3\mu a_0 U, \quad e_3 - e_1 = -3\mu a_0 V, \quad e_1 - e_2 = -3\mu a_0 W
\]
und daraus wegen $e_1+e_2+e_3=0$,
\[
\begin{aligned}
e_1 &= \mu a_0 (V-W),\\
e_2 &= \mu a_0 (W-U),\\
e_3 &= \mu a_0 (U-V).
\end{aligned}
\]
